\chapter{Gastric Bypass Surgery}

\section*{Overview}
Roux-en-Y gastric bypass is metabolic and bariatric surgery that creates a small stomach pouch and reroutes part of the small intestine.

\section*{Common Instruments and Equipment}
The exact setup varies by technique and setting. Common items may include:
\begin{itemize}[nosep,leftmargin=*]
  \item laparoscopic or open bariatric surgery instruments
  \item surgical staplers and bowel-clamping instruments
  \item bougie or sizing tube and leak-test equipment when used
\end{itemize}

\section*{Why It Is Done}
It may be considered for severe obesity with related disease when structured medical and lifestyle treatment has not achieved adequate or durable benefit.

\section*{How It Is Performed}
Usually laparoscopically, the surgeon divides the stomach to create a small pouch and connects it to a lower segment of small intestine. Food bypasses most of the stomach and the first part of the intestine.

\section*{Contraindications}
No single contraindication applies to every patient; the treating team assesses whether the procedure is suitable for the individual and the reason for treatment.

\section*{Precautions}
The care team reviews relevant health conditions, medicines, allergies, and procedure-specific risks, and explains any changes needed for the individual.

\section*{Risks and Recovery}
Bleeding, infection, leaks, clots, obstruction, dumping syndrome, ulcers, gallstones, and nutritional deficiencies are possible. Lifelong dietary guidance, vitamin and mineral supplementation, laboratory monitoring, and follow-up are required.

\section*{References}
\begin{enumerate}
  \item MedlinePlus. Gastric Bypass Surgery. U.S. National Library of Medicine; 2025.
  \item Current Medical Diagnosis and Treatment. McGraw-Hill; 2025.
\end{enumerate}
\clearpage
